\chapter{Baseline Architectures and Implementation}
\label{chap:architecture}

This chapter describes the two baseline execution strategies in full:
their architecture, their pseudocode, the exact prompts issued to the
LLM, and the implementation layer (SQL parsing, the multi-modal operator
abstraction, and the evaluation harness) shared by both. Unlike earlier
drafts of this report, every prompt and every architectural claim in
this chapter is taken verbatim or paraphrased directly from the project
repository, branch \texttt{agent/consolidate-four-method-benchmarks}:
\url{https://github.com/AnnaRemi/ENSIMAG-AI-M2-final-project/tree/agent/consolidate-four-method-benchmarks}.
Where an earlier draft of this report guessed at an implementation
detail that turned out to differ from the real code (most notably the
Trummer baseline's join structure, corrected in
Section~\ref{sec:trummer-baseline}), the correction is noted explicitly
rather than silently.

\section{System Overview}

Figure~\ref{fig:system-overview} shows the overall system: a natural
language query is decomposed into a structured predicate $\sigma_S$ and a
semantic predicate $\sigma_U$ (predicate extraction), the structured part
is compiled to SQL and run first, and the semantic part is then evaluated
by one of the four strategies described in this chapter and the next.

\begin{figure}[h]
\centering
\resizebox{\textwidth}{!}{%
\begin{tikzpicture}[
  node distance=7mm and 8mm,
  box/.style={draw, rounded corners, align=center, minimum height=9mm, font=\small, fill=blue!6},
  llmbox/.style={draw, rounded corners, align=center, minimum height=9mm, font=\small, fill=orange!12},
  arr/.style={-{Latex[length=2.2mm]}}
]
\node[box] (nl) {Natural-language\\query $Q$};
\node[llmbox, right=of nl] (extract) {Predicate\\extraction\\$Q \to (\sigma_S,\sigma_U)$};
\node[box, right=of extract] (planner) {Access planner /\\query compiler};
\node[box, right=of planner] (sql) {SQL predicate\\pushdown\\($R \to R_{\sigma_S}$)};
\node[box, below=14mm of sql] (strategy) {Semantic operator\\(SUQL / Trummer /\\cascaded variant)};
\node[box, right=of strategy] (out) {JSON\\result set};
\draw[arr] (nl)--(extract)--(planner)--(sql);
\draw[arr] (sql) -- (strategy);
\draw[arr] (strategy) -- (out);
\end{tikzpicture}}
\caption{End-to-end system: NL query $\to$ predicate extraction $\to$ SQL
pushdown $\to$ pluggable semantic-operator strategy $\to$ JSON output.
The four strategies compared in this report (SUQL baseline/V1, Trummer
baseline/V1) differ in the ``semantic operator'' box, and --
importantly, see Section~\ref{sec:trummer-baseline} -- in whether the
SQL-pushdown box is applied at all before that operator runs.}
\label{fig:system-overview}
\end{figure}

For SUQL (baseline and V1), predicate extraction is itself an LLM call
(\texttt{nl\_to\_suql}, a few-shot semantic parser that compiles the NL
question directly into a SUQL query -- plain SQL \texttt{WHERE} clauses
for structured columns, plus \texttt{answer()}/\texttt{summary()}
free-text predicates), and the parser's own system prompt explicitly
instructs the model to place structured predicates before \texttt{answer()}
in the \texttt{WHERE} clause ``for efficiency.'' The structural SQL is
then stripped out and executed on an in-memory SQLite instance
(``litesql'') \emph{before} any \texttt{answer()} call is issued --
Section~\ref{sec:suql-baseline} gives the exact pipeline. For Trummer
(baseline and V1), predicate extraction and structural filtering are
handled differently, and, critically, \emph{not at all} in the baseline
-- see Section~\ref{sec:trummer-baseline}.

\section{SUQL Baseline: Row-Wise Predicate Pushdown}
\label{sec:suql-baseline}

\begin{figure}[h]
\centering
\resizebox{\textwidth}{!}{%
\begin{tikzpicture}[
  node distance=6mm and 6mm,
  box/.style={draw, rounded corners, align=center, minimum height=8mm, font=\small, fill=blue!6},
  llmbox/.style={draw, rounded corners, align=center, minimum height=8mm, font=\small, fill=orange!12},
  arr/.style={-{Latex[length=2mm]}}
]
\node[box] (q) {$R_{\sigma_S}$\\(post pushdown)};
\node[box, right=of q] (join) {Row-wise join\\with $D$ on \textit{id}};
\node[llmbox, right=of join] (udf) {\answerfn{}(review,\\$\sigma_U$) -- one\\call per row};
\node[box, right=of udf] (out) {Yes/No\\per row};
\node[box, right=of out] (filt) {Keep rows\\where Yes};
\draw[arr] (q)--(join)--(udf)--(out)--(filt);
\end{tikzpicture}}
\caption{SUQL baseline: after SQL pushdown, every surviving row is joined
to its review and judged by one \answerfn{} call.}
\label{fig:suql-baseline}
\end{figure}

\textbf{Pipeline.} The engine (\texttt{project SUQL/baseline/suql\_engine.py})
implements the SUQL paper's own pipeline in five steps: (1) split the
compiled SUQL query into a structural SQL part and a set of
\texttt{answer()}/\texttt{summary()} free-text predicates; (2) run the
structural SQL on an in-memory SQLite table; (3) for every row that
survives, evaluate each \texttt{answer()} predicate with one LLM call,
short-circuiting the remaining predicates for a row the moment one of
them returns No; (4) keep only rows that pass every \texttt{answer()}
predicate; (5) compute \texttt{summary()} values for any requested
summary columns. Row-level results are cached in-process on
\texttt{(review\_hash, question)} to avoid re-querying identical pairs.

\textbf{Cost model.} For a single \texttt{answer()} predicate, the
number of semantic-filter LLM calls is exactly $|R_{\sigma_S}|$, one per
structurally-surviving row; wall-clock time is this many calls executed
serially against a single Ollama-served model instance. In practice the
harness's \texttt{llm\_calls} metric (Section~\ref{sec:eval-harness})
counts \emph{every} LLM prompt the pipeline issues for a question --
including the one-time semantic-parsing call and any \texttt{summary()}
calls -- so the measured call count is a modest constant multiple of
$|R_{\sigma_S}|$ rather than exactly $|R_{\sigma_S}|$; Chapter~\ref{chap:experiments}
reports both the raw measured totals and, in
Section~\ref{sec:theory-cost}, a ``lean'' one-call-per-row idealization
used for the theoretical cost analysis.

\textbf{Quality.} Because every row is judged individually, with the
full review text in context and no competing rows to distract the model,
this strategy is close to the achievable quality ceiling for a given
base model; empirically (Chapter~\ref{chap:experiments}) it is a solid,
consistent mid-table performer among the four studied variants.

\begin{algorithm}[h]
\caption{SUQL baseline: row-wise predicate pushdown}
\label{alg:suql-baseline}
\begin{algorithmic}[1]
\Require structured relation $R$, review relation $D$, structured
predicate $\sigma_S$, semantic predicate $\sigma_U$
\State $R_{\sigma_S} \gets \textsc{ExecuteSQL}(\texttt{SELECT * FROM } R \texttt{ WHERE } \sigma_S)$
\State $\mathit{result} \gets \emptyset$
\For{$r \in R_{\sigma_S}$}
  \State $\mathit{review} \gets D[\,r.\mathit{id}\,]$ \Comment{join on id}
  \If{$(\mathit{review}, \sigma_U) \in \mathit{cache}$}
    \State $b \gets \mathit{cache}[(\mathit{review}, \sigma_U)]$
  \Else
    \State $b \gets \textsc{LLMCall}\bigl(\texttt{answer}(\mathit{review}, \sigma_U)\bigr)$ \Comment{1 LLM call}
    \State $\mathit{cache}[(\mathit{review}, \sigma_U)] \gets b$
  \EndIf
  \If{$b = \textsc{Yes}$}
    \State $\mathit{result} \gets \mathit{result} \cup \{r\}$
  \EndIf
\EndFor
\State \Return $\mathit{result}$
\end{algorithmic}
\end{algorithm}

\begin{lstlisting}[style=prompt, caption={\answerfn{} system prompt, verbatim from \texttt{project SUQL/baseline/suql\_engine.py} (SUQL baseline).}, label={lst:suql-prompt}]
You are a recall-first semantic filter for movie reviews.

Decide whether the review provides evidence that the movie
satisfies the question.

Recall is the primary objective, but do not accept every
review:
- Answer YES when there is any concrete review-specific
  evidence: direct wording, a synonym/paraphrase, a described
  example, or a reasonable implication.
- Give borderline evidence the benefit of the doubt when it
  specifically bears on the condition. The evidence need not
  use the question's exact words.
- This is evidence retrieval, not sentiment scoring: an
  explicit mention still counts when it is negated, qualified,
  quoted, or used critically.
- Answer NO when there is no review-specific support, the
  connection is only the movie's genre/topic, the text is
  generic praise or criticism, or the review never discusses
  the condition itself.
- Do not infer evidence merely because the condition is not
  contradicted.

Return exactly YES or NO. Do not include reasoning,
punctuation, or extra text.
\end{lstlisting}

Two design choices in this prompt matter for the cost/quality
discussion in later chapters. First, it is explicitly
\emph{recall-first}: it instructs the model to give ambiguous evidence
the benefit of the doubt, which biases the baseline toward higher recall
at some precision cost. Second, it asks for a bare \texttt{YES}/\texttt{NO}
token rather than a JSON object or a justification string, which keeps
the completion short (a handful of output tokens) and avoids the
JSON-parsing failure mode that a richer response schema would introduce.

\section{Trummer Baseline: Adaptive Two-Collection Block Join}
\label{sec:trummer-baseline}

\textbf{Correction relative to earlier drafts.} An earlier draft of this
report described the Trummer baseline as operating on the
\emph{structurally pre-filtered} relation $R_{\sigma_S}$, serialized into
a single header-plus-rows chunk. This is incorrect for the actual
implementation (\texttt{project Trummer/baseline/trummer\_join/}), and
the correction matters: the Trummer baseline is, by design, the
\emph{no-pushdown} reference point in this study. Its own README states
plainly that it uses ``no SUQL-style structured pruning or cascade''; the
structured relation and the review relation are kept as two
\emph{separate} logical collections, and the LLM itself is asked to
perform the join, mirroring the two-source, multi-modal retrieval
setting that ELEET's operator abstraction targets
(Section~\ref{sec:formal-model}).

\begin{figure}[h]
\centering
\resizebox{\textwidth}{!}{%
\begin{tikzpicture}[
  node distance=6mm and 6mm,
  box/.style={draw, rounded corners, align=center, minimum height=8mm, font=\small, fill=blue!6},
  llmbox/.style={draw, rounded corners, align=center, minimum height=8mm, font=\small, fill=orange!12},
  arr/.style={-{Latex[length=2mm]}}
]
\node[box] (r) {Collection 1:\\structured movie\\rows (block 1)};
\node[box, below=8mm of r] (d) {Collection 2:\\review rows\\(block 2)};
\node[llmbox, right=16mm of r] (join) {LLM block-pair join:\\match \texttt{tconst}=\texttt{movie\_id},\\apply $\sigma_U$, emit\\\texttt{movie\_id: YES/NO}};
\node[box, right=of join] (parse) {Line-oriented\\parse, union\\block pairs};
\draw[arr] (r) -- (join);
\draw[arr] (d) -- (join);
\draw[arr] (join) -- (parse);
\end{tikzpicture}}
\caption{Trummer baseline: no structural pushdown. The full candidate
pool is split into blocks of each collection independently; the LLM is
asked to join a block of Collection~1 against a block of Collection~2
\emph{and} apply the semantic predicate in the same call.}
\label{fig:trummer-baseline}
\end{figure}

\textbf{Pipeline.} The candidate movies (Collection~1) and their reviews
(Collection~2) are each partitioned into blocks via
\texttt{partition\_records}, with block size chosen by an
\texttt{optimal\_block\_size} routine (a port of an existing
\texttt{llmjoin} block-size tuning utility). \texttt{join\_two\_blocks}
sends one block of each collection to the LLM in a single prompt
(Listing~\ref{lst:trummer-prompt}) and parses a line-oriented
\texttt{movie\_id: YES} / \texttt{movie\_id: NO} response with a
tolerant regex-based parser (\texttt{parse\_pairs}/\texttt{parse\_decisions}),
so execution does not depend on provider-specific structured-output
support. \texttt{adaptive\_join} (rather than the simpler
\texttt{block\_join}) additionally checks response coverage
(\texttt{has\_complete\_decision\_coverage}) and adjusts block pairing
when a response fails to decide every pair it was asked about -- the
``adaptive'' half of the name.

\textbf{Cost model.} Because there is no structural pushdown, the number
of LLM calls scales with the number of block \emph{pairs} needed to
cover the \emph{entire} candidate pool (up to 100 rows in the 10-query
benchmark of Chapter~\ref{chap:experiments}, not just the 40 that
satisfy $\sigma_S$), and grows further whenever a block pair's response
fails coverage and must be retried or re-partitioned -- this is the main
reason the Trummer baseline's call count is both higher, on average,
than a naive $|R|/(\text{block size})$ estimate would suggest, and
highly variable across questions (Chapter~\ref{chap:experiments} reports
a range from $30$ to $268$ calls across the ten held-out questions).

\textbf{Quality.} Two systematic error sources appear, both traceable to
the fact that the model must recover row identity and column structure
from block text and perform the id-matching join itself, without any
guarantee of covering every pair in one pass:
\begin{enumerate}
  \item \textbf{Join/parsing errors.} A block pair may return
        \texttt{movie\_id}s that do not exactly match any row in
        Collection~1 (drift in how the model echoes the id), which the
        tolerant parser discards rather than guesses.
  \item \textbf{Incomplete coverage.} The model may decide only a subset
        of the pairs in a block before truncating or drifting, which --
        absent structural pushdown to shrink the pool the join has to
        cover in the first place -- is the dominant driver of the
        Trummer baseline's low recall (Chapter~\ref{chap:experiments}).
\end{enumerate}

\begin{algorithm}[h]
\caption{Trummer baseline: adaptive two-collection block join}
\label{alg:trummer-baseline}
\begin{algorithmic}[1]
\Require structured relation $R$ (\emph{not} pre-filtered by $\sigma_S$
inside this operator), review relation $D$, semantic predicate $\sigma_U$
\State $B_1 \gets \textsc{PartitionRecords}(R, \textsc{OptimalBlockSize}(R))$
\State $B_2 \gets \textsc{PartitionRecords}(D, \textsc{OptimalBlockSize}(D))$
\State $\mathit{result} \gets \emptyset$
\For{$(b_1, b_2) \in B_1 \times B_2$ \textbf{as scheduled by} \textsc{AdaptiveJoin}}
  \State $\mathit{resp} \gets \textsc{LLMCall}\bigl(\text{join-prompt}(b_1, b_2, \sigma_U)\bigr)$ \Comment{1 LLM call per block pair}
  \If{$\lnot \textsc{HasCompleteCoverage}(\mathit{resp}, b_1)$}
    \State \textsc{AdjustBlockPairing}() \Comment{retry/re-partition}
  \EndIf
  \State $\mathit{pairs} \gets \textsc{ParsePairs}(\mathit{resp}, b_1, b_2)$
  \State $\mathit{result} \gets \mathit{result} \cup \{\, r_1 : (r_1,r_2) \in \mathit{pairs} \,\}$
\EndFor
\State \Return $\mathit{result} \cap R_{\sigma_S}$ \Comment{$\sigma_S$ applied post hoc, if at all, by the harness}
\end{algorithmic}
\end{algorithm}

\begin{lstlisting}[style=prompt, caption={Block-join prompt, verbatim from \texttt{project Trummer/baseline/trummer\_join/operators.py} (\texttt{create\_prompt}).}, label={lst:trummer-prompt}]
Find every movie in Collection 1 that has a review in
Collection 2 satisfying this predicate:
{predicate}
Act as a recall-first semantic filter, while requiring
review-specific evidence.
- Answer YES for direct wording, synonyms/paraphrases,
  described examples, or reasonable implications.
- Give borderline evidence the benefit of the doubt when it
  specifically bears on the predicate.
- Count explicit mentions even when negated, qualified,
  quoted, or critical; retrieve evidence rather than score
  sentiment.
- Answer NO when support is absent, generic, based only on
  genre/topic, or never discusses the predicate itself.
- Lack of contradiction alone is not evidence.
Collection 1 contains structured movie rows. Collection 2
contains review chunks.
A review belongs to a movie only when tconst is exactly
equal to movie_id.
Return one decision per line as: movie_id: YES or
movie_id: NO.
Do not include confidence, evidence, explanations, or prose.
Copy each movie_id exactly from Collection 1.
Do not add prose, JSON, markdown, or a completion marker.

Collection 1:
{block_1_rows}

Collection 2:
{block_2_rows}

Decisions:
\end{lstlisting}

Note the shared evidence-retrieval language between this prompt and the
SUQL baseline's \texttt{answer()} prompt (Listing~\ref{lst:suql-prompt}):
both explicitly frame the task as recall-first evidence retrieval rather
than sentiment scoring, and both instruct the model to count negated,
qualified, or critical mentions as evidence. This is a deliberate,
shared design choice across both projects in the repository, and it
means the quality gap between the two baselines (Chapter~\ref{chap:experiments})
is attributable to the \emph{execution strategy} -- row-wise judgment
with pushdown vs.\ block-wise LLM-driven join without pushdown -- and
not to a difference in how ``satisfies the predicate'' is defined for
the underlying model.

\section{Shared Evaluation Harness}
\label{sec:eval-harness}

All four variants (both baselines here, and both cascaded V1 variants of
Chapter~\ref{chap:cascading}) are evaluated with the same harness
(\texttt{benchmarks/shared/scripts/}):
\begin{itemize}
  \item \textbf{Gold set derivation.} Rather than hardcoding an expected
        positive count (an early version of the evaluation script fixed
        \texttt{actual\_p = 11} for a specific query and had to be
        corrected), the current benchmark suites derive ground truth by
        construction: every question's candidate pool is assembled from
        annotated rows whose \texttt{structured\_match} and
        \texttt{semantic\_label} columns are known in advance
        (Section~\ref{sec:10q-suite}), and \texttt{ground\_truth} is
        their logical conjunction. This keeps the confusion-matrix
        computation (TP, FP, FN, TN) parameterized by the actual
        question and its annotated pool, rather than by a magic
        constant, and additionally makes every label auditable via an
        \texttt{evidence\_excerpt} and an \texttt{annotation\_rationale}.
  \item \textbf{Metrics.} Precision, recall, F1 against the gold set;
        wall-clock time per question; total LLM calls per question, plus
        (for the cascaded variants) the cheap/expensive split of both
        calls and wall time -- computed identically for all four
        variants and averaged over repeated runs
        (Section~\ref{sec:10q-suite}).
  \item \textbf{Implementation isolation.} Each method's engine module is
        imported fresh per run with its own environment configuration
        (model names, API base, cascade parameters), and its in-process
        cache is cleared before each run, so repeated executions are not
        silently accelerated by stale cached answers from a previous
        method or question.
\end{itemize}
